% !TEX root = ../bb_AiRa_Kappaleet.tex

%% HARRIS: Chapter 10

% ö = \%c3\%b6
% ä = \%C3\%A4

\xdef\sectiontitle{\IIsectiontitle} 
\TOCrefs

%%%%%%%%%%%%%%%
\begin{frame}[label=2_8_a]%[label=2_7_TOC1]
\frametitle{\texorpdfstring{\culine[\sectiontitleulinecolor]{\sectiontitle}}{\sectiontitle} }
\label{\TOCname}

\vspace{\chapterTOCvksip}
\begin{itemize}[leftmargin=0.5cm,itemsep=3mm]
    \item[2.1]<1-> \hyperlink{2_1_a}{\IIaaSubsectiontitle}
    \item[2.2]<1-> \hyperlink{2_2_a}{\IIaSubsectiontitle}
    \item[2.3]<1-> \hyperlink{2_3_a}{\IIbSubsectiontitle}	
    \item[2.4]<1-> \hyperlink{2_4_a}{\IIcSubsectiontitle}	
    \item[2.5]<1-> \hyperlink{2_5_a}{\IIdSubsectiontitle}
    \item[2.6]<1-> \hyperlink{2_6_a}{\IIeSubsectiontitle}
    \item[2.7]<1-> \hyperlink{2_7_a}{\IIfSubsectiontitle}
    \item[2.8]<1-> \hyperlink{2_8_a}{\CDAlert[blue]{\IIgSubsectiontitle}}
    \item[2.9]<1-> \hyperlink{2_9_a}{\IIhSubsectiontitle}
    \item[2.10]<1-> \hyperlink{2_10_a}{\IIiSubsectiontitle}
\end{itemize}


%\endofslide 
\end{frame}
%%%%%%%%%%%%%%%

\xdef\sectiontitle{\IIgSubsectiontitle} 

%%%%%%%%%%%%%%%
\begin{frame}
\frametitle{\texorpdfstring{\culine[\sectiontitleulinecolor]{\sectiontitle}}{\sectiontitle} }
\framesubtitle{Binääriset yhdistepuolijohteet}
\appear{}
%Compound semiconductors
\begin{itemize}[itemsep=1.5mm] \seminormal
	\item Alkuainepuolijohteet kuuluvat suurimmaksi osaksi jaksollisen järjestelmän ryhmään 14 (IV): \appear{Si, Ge}
	
\item \CDAlert[fuchsia]{Yhdistepuolijohteet}, joista suurin osa on ryhmien III–V\appear{, II–VI} \appear{tai IV–IV alkuaineiden sekoituksia}\appear{, ovat erittäin yleisiä} \appear{(myös Si$_{1-x}$Ge$_{x}$)}
\appear{\xdef\loc{south east}%
\begin{tikzpicture}[remember picture,overlay] % anchor=south
    \node (A) [yshift=-0.5*\figYshift,xshift=0mm, anchor=\loc] at (current page.\loc) {\includegraphics[page=1,height=5.65cm]{Figures_booklet/SOM_semiconductor_table_figure.png}};%scale=\figscale. % 8cm max height
\end{tikzpicture}}
\item Muuttamalla yhdisteiden alkuainepitoisuuksia, puolijohteen ominaisuuksiin voidaan suuresti vaikuttaa \\ 
\end{itemize}
% 6.5 cm = midway, 10 cm = 3/4 
\vspace{1.3mm}
\begin{minipage}{5.35cm}
	\begin{itemize}[itemsep=1.5mm]
	\item[] (erityisesti energiarakoon\appear{, ja vöiden kaareutumisiin)}

\item Kuvassa nähdään \Alert[red]{energiarako \appear{hilavakion funktiona}} \appear{yleisimmille binäärisille yhdistepuolijohteille\footnote{www.tf.uni-kiel.de} }

\item Ominaisuuksien säädettävyys on oleellista puolijohdekomponenttien valmistuksessa!
\end{itemize}

\end{minipage}

% great WWW: https://www.tf.uni-kiel.de/matwis/amat/semi_en/kap_5/backbone/r5_1_4.html
%\endofslide 
%\endofsection
\end{frame}
%%%%%%%%%%%%%%%

%%%%%%%%%%%%%%%
\begin{frame}
\frametitle{\texorpdfstring{\culine[\sectiontitleulinecolor]{\sectiontitle}}{\sectiontitle} }
\framesubtitle{Diodi-liitos}
%The Diode
 \appear{}
 \seminormal
\begin{itemize}
	\item \CDAlert[fuchsia]{Diodi} \appear{(tai diodi-liitos)} \appear{on ehkäpä yleisin puolijohdekomponentti.} \appear{Sen läpi sähkövirta \Alert[fuchsia]{kulkee vain yhteen suuntaan} }
	
	\item Diodi rakentuu p- \appear{ja n-tyypin puolijohdeosioista}\appear{, jotka on liitetty atomien tasolla kiinni toisiinsa.}\appear{ Yhdessä p- ja n-tyypin puolijohteista muodostuu \CDAlert[fuchsia]{pn-liitos}}	
	
	\item \CDAlert[red]{Päästösuunnattuun} pn-liitokseen kohdistetaan jännite \appear{niin että p-tyypin osan potentiaali on korkeammalla:}
\end{itemize}
	
% 6.5 cm = midway, 10 cm = 3/4 
\vspace{2.5mm}
\begin{minipage}{6.3cm}
	\begin{itemize}
	\item[] \begin{itemize}	
	\item elektronit sekä aukot mo-\\
		lemmat etenevät kohti liitos-\\
		kohtaa
	\item Kohdatessaan ne yhdistyvät 
	\implies{ \Alert[red]{liitoksen yli muodostuu \\jatkuva virta } } 
\end{itemize}
\end{itemize}

\end{minipage}

\xdef\loc{south east}
\begin{tikzpicture}[remember picture,overlay] % anchor=south
    \node (A) [yshift=-0.25*\figYshift,xshift=\figXshift, anchor=\loc] at (current page.\loc) {\includegraphics[page=1,width=8cm]{Figures_booklet/SOM_10_Bonding_figures/SOM_Bonding_Figure_41.png}}; %scale=\figscale. % 8cm max height
\end{tikzpicture}
\footnote{\HarrisCR[1-]}

\endofslide 
%\endofsection
\end{frame}
%%%%%%%%%%%%%%%

%%%%%%%%%%%%%%%
\begin{frame}
\frametitle{\texorpdfstring{\culine[\sectiontitleulinecolor]{\sectiontitle}}{\sectiontitle} }
\framesubtitle{Diodi-liitos}
%The Diode
 
 \seminormal
\begin{itemize}
	\item<1-> \CDAlert[fuchsia]{Diodi} (tai diodi-liitos) on ehkäpä yleisin puolijohdekomponentti. Sen läpi sähkövirta \Alert[fuchsia]{kulkee vain yhteen suuntaan}
	
	\item<1-> Diodi rakentuu p-tyypin puolijohdepalasta sekä n-tyypin puolijohteesta, jotka on liitetty atomien tasolla kiinni toisiinsa. Yhdessä p- ja n-tyypin puolijohteista muodostuu \CDAlert[fuchsia]{pn-liitos}	
		
	\appear{}
	\item Pn-liitos voidaan kytkeä \appear{myös niin että $n$-tyypin puoli on korkeammassa potentiaalissa}\appear{, jolloin liitos on \CDAlert[blue]{estosuunnattu} }
	\implies{ Sekä elektronit että aukot liik- \\
		      kuvat liitoksesta poispäin, ja \\
		      \Alert[blue]{virrankulku pysähtyy lähes heti}  }
		      
	\item Diodin käytös on todella erilaista \\ 
		kuin esim. vastuksien\appear{, jotka\\
		noudattavat Ohmin lakia}
\end{itemize}
	
\xdef\loc{south east}
\begin{tikzpicture}[remember picture,overlay] % anchor=south
    \node (A) [yshift=-0.25*\figYshift,xshift=\figXshift, anchor=\loc] at (current page.\loc) {\includegraphics[page=1,width=8cm]{Figures_booklet/SOM_10_Bonding_figures/SOM_Bonding_Figure_41.png}}; %scale=\figscale. % 8cm max height
\end{tikzpicture}
\footnote{\HarrisCR[1-]}

\endofslide 
%\endofsection
\end{frame}
%%%%%%%%%%%%%%%

%%%%%%%%%%%%%%%
\begin{frame}
\frametitle{\texorpdfstring{\culine[\sectiontitleulinecolor]{\sectiontitle}}{\sectiontitle} }
\framesubtitle{pn-liitos}
%p-n junction
\appear{}
\seminormal
\begin{itemize}[itemsep=1.4mm]
	\item Kaksi seostettua puolijohdetta voidaan liittää toisiinsa \appear{\CDAlert[fuchsia]{synnyttäen} \CDAlert[fuchsia]{pn-liitoksen}} \appear{(esim. p-tyypin Ge ja n-tyypin Ge)}
	
	\item Erillään\appear{ kappaleilla on toisistaan poikkeavat fermitasot $E_F$} 
	\appear{\xdef\loc{south east}%
\begin{tikzpicture}[remember picture,overlay] % anchor=south
    \node (A) [yshift=-0.35*\figYshift,xshift=1*\figXshift, anchor=\loc] at (current page.\loc) {\includegraphics[page=1,height=5.8cm]{Figures_booklet/SOM_10_Bonding_figures/SOM_Bonding_Figure_42.png}}; %scale=\figscale. % 8cm max height
\end{tikzpicture}\footnote{\HarrisCR}}%

	\item Liitoksen jälkeen, n-tyypissä korkealla olevat elektronit siirtyvät p-tyypin puolelle ja vastaavasti aukot
\end{itemize}

% 6.5 cm = midway, 10 cm = 3/4 
%\vspace{2.5mm}
\begin{minipage}{5.4cm}
	\begin{itemize}[itemsep=1.4mm]
	\item[] siirtyvät p-tyypin puolelta \\
		   n-tyypin puolelle 
	
	\item Lyhyen varauksenkuljettajien siir\-ty\-mähetken jälkeen\appear{ systeemi löytää uuden makroskooppisen tasapainotilan}
	
	\item Tasapainossa fermitasot ovat molemmilla puolilla liitosta samat \appear{muodostaen kuvan (a) energiatasokaavion}
	
%	\item In addition, there will be a small build up of negative and positive charges into either sides of the depletion zone
\end{itemize}
\end{minipage}

\endofslide 
%\endofsection
\end{frame}
%%%%%%%%%%%%%%%

%%%%%%%%%%%%%%%
\begin{frame}
\frametitle{\texorpdfstring{\culine[\sectiontitleulinecolor]{\sectiontitle}}{\sectiontitle} }
\framesubtitle{Estosuuntainen pn-liitos}
\appear{}
\begin{itemize}[itemsep=1.5mm]
	\item Kuvan 42 energiatasokaaviot selventävät diodin toimintaa. 
	 \appear{Kytketyn nollajännitteen tapauksessa (a)}\appear{ vain muutamien elektronien täytyy siirtyä liitoksen yli, jotta tasapainotila saavutetaan}
	
	\item Liitosaluetta kutsutaan \Alert[red]{tyhjennysalueeksi}, jossa ei ole  varauksen-\\
\end{itemize}

\begin{minipage}{5.3cm}
	\begin{itemize}[itemsep=1.5mm]
	\item[] kuljettajia \appear{(leveys estää elektronien tunne\-loi\-tu\-misen)}
	%, vöiden eroja on kuvassa liioiteltu
	\item \Alert[blue]{Estosuunnatussa liitok\-ses\-sa} (b) \appear{vähemmistöva\-rauk\-sen\-kuljettajat ovat har\-vi\-naisia}\appear{, joten eivät synnytä virtaa.} \appear{Kasvanut potentiaaliero puolestaan estää enemmistövarauksenkuljettajien liikkeen}

\end{itemize}
\end{minipage}

\xdef\loc{south east}%
\begin{tikzpicture}[remember picture,overlay] % anchor=south
    \node (A) [yshift=-0.35*\figYshift,xshift=\figXshift, anchor=\loc] at (current page.\loc) {\includegraphics[page=1,height=5.8cm]{Figures_booklet/SOM_10_Bonding_figures/SOM_Bonding_Figure_42.png}};%scale=\figscale. % 8cm max height
\end{tikzpicture}
\footnote{\HarrisCR[1-]}
\endofslide 
%\endofsection
\end{frame}
%%%%%%%%%%%%%%%

%%%%%%%%%%%%%%%
\begin{frame}
\frametitle{\texorpdfstring{\culine[\sectiontitleulinecolor]{\sectiontitle}}{\sectiontitle} }
\framesubtitle{Päästösuuntainen pn-liitos}
%\appear{}
\begin{itemize}[itemsep=1.5mm]
	\item<1-> Kuvan 42 energiatasokaaviot selventävät diodin toimintaa. Kytketyn nollajännitteen tapauksessa (a) vain muutamien elektronien täytyy siirtyä liitoksen yli, jotta tasapainotila saavutetaan
	
	\item<1-> Liitoksen kutsutaan called \Alert[red]{tyhjennysalueeksi}, jossa ei ole  varauksen-\\
\end{itemize}
%\vspace{2.5mm}
\begin{minipage}{5.3cm}
	\begin{itemize}[itemsep=1.5mm]
	\item<1->[] kuljettajia (leveys estää elektronien tunne\-loi\-tu\-misen)
	
%	\item \Alert[blue]{Estosuuntaisella liitoksella} (b, vöiden eroja on kuvassa liioiteltu)\appear{ vähemmistövarauksenkuljettajat ovat harvinaisia}

	\appear{}
	\item \Alert[red]{Päästösuunnatussa liitoksessa} (c) n-tyypin osion johtavuusvyö nousee ylöspäin  %relative to that of the $\mathrm{p}$-type 
	\implies{ Elektronit ja aukot \CDAlert[red]{tyhjen-}\\ 
		      \CDAlert[red]{nysalueen lähellä} \\
		      \appear{diffundoituvat sen yli}\appear{, \\%diffuse cross
		      yhdistyvät ja synnyttävät \\
		      \CDAlert[red]{virran} }}
\end{itemize}
\end{minipage}

\xdef\loc{south east}
\begin{tikzpicture}[remember picture,overlay] % anchor=south
    \node (A) [yshift=-0.35*\figYshift,xshift=\figXshift, anchor=\loc] at (current page.\loc) {\includegraphics[page=1,height=5.8cm]{Figures_booklet/SOM_10_Bonding_figures/SOM_Bonding_Figure_42.png}}; %scale=\figscale. % 8cm max height
\end{tikzpicture}
\footnote{\HarrisCR[1-]}

\endofslide 
%\endofsection
\end{frame}
%%%%%%%%%%%%%%%

%%%%%%%%%%%%%%%
\begin{frame}
\frametitle{\texorpdfstring{\culine[\sectiontitleulinecolor]{\sectiontitle}}{\sectiontitle} }
\framesubtitle{Sähkövirrat tyhjennysalueella}%of the p-n junction
\appear{}
\seminormal%
\begin{itemize}[itemsep=1.5mm]
	\item pn-liitoksessa on itseasiassa neljä erilaista sähkövirtaa: \appear{kahdentyyppisiä virtoja} \appear{(ajelehtimisvirtoja ja diffuuseja)} \appear{sekä \CDAlert[red]{enemmistövarauksenkuljettajille}} \appear{että \CDAlert[blue]{vähemmistövarauksenkuljettajille}} 
	\appear{\xdef\loc{south east}%
\begin{tikzpicture}[remember picture,overlay] % anchor=south
    \node (A) [yshift=-0.25*\figYshift,xshift=1*\figXshift, anchor=\loc] at (current page.\loc) {\includegraphics[page=2,width=5cm]{Figures_booklet/Tikz_SOM_Bonding_figures_pn_junction}};%
\end{tikzpicture}%
}

	\item Elektronit poikkeavat \appear{ulkoiseen kenttään nähden vastakkaiseen suuntaan} \appear{($F=qE$)} \appear{n-puolella}\appear{, ja aukot poikkeavat samansuuntaisesti kentän kanssa}
	
	\item Tilojen energia\appear{/konsentraatioeroista} \appear{johtuva diffuusio} \appear{johtaa \CDAlert[red]{rekombinaatiovirtoihin}} \appear{($e^-$:lle ja $h^+$:lle)}\appear{, joita merkitään termeillä $i_{\text {aukot,rek}} \Equiv i_{p, r}$} \appear{ja $i_{\text {elektronit,rek}} \Equiv i_{n, r}$}

\item Elektroni–aukko-pareja syntyy myös liitoksen \\
	lähellä termisten virittymisten johdosta%\appear{, and are swept through the junction due to the build-up E-field in the \CDAlert[fuchsia]{depletion region} }

 
\end{itemize}

\vspace{1.5mm}
\begin{minipage}{8.45cm}
	\begin{itemize}[itemsep=1.5mm]
\item Näitä virtoja kutsutaan \appear{\Alert[blue]{(termisiksi)} \CDAlert[blue]{ge\-ne\-raa-} \CDAlert[blue]{tiovirroiksi}, $i_{\text {p,gen}} \Equiv i_{p,g} \quad$} \appear{ja $\quad i_{\text {e,gen}} \Equiv i_{n,g}$}

\item Tasapainossa\appear{ generaatio- ja rekombinaatiovirrat ovat yhtäsuuria:} %
$
 \appear{| i_{p, g} |}  \equal \appear{| i_{p, r} |} \appear{~ja~} \appear{| i_{n, g} |}  \equal  \appear{| i_{n, r} |}
$
\end{itemize}
\end{minipage}


\endofslide 
%\endofsection
\end{frame}
%%%%%%%%%%%%%%%

%% LEAVE DIFFUSION REGIONS AND THE DEPLETION REGION TO FOLLOW-UP COURSES??
%%%%%%%%%%%%%%%%
%\begin{frame}
%\frametitle{\texorpdfstring{\culine[\sectiontitleulinecolor]{\sectiontitle}}{\sectiontitle} }
%
%\begin{itemize}
%	\item Taking a closer look at the junction area, There are actually three zones:
%
%\item Depletion region in the middle, across the junction
%
%\item Diffusion regions on both sides of the depletion region, from there the electrons and holes diffuse to the junction area
%
%\item Homogeneous regions far away from Home the junction, which are actually not much affected by the junction
%\end{itemize}
%
%%\xdef\loc{south east}
%%\begin{tikzpicture}[remember picture,overlay] % anchor=south
%%    \node (A) [yshift=-0.25*\figYshift,xshift=1*\figXshift, anchor=\loc] at (current page.\loc) {\includegraphics[page=6,height=4cm]{Figures_booklet/Tikz_SOM_Bonding_figures}}; %scale=\figscale. % 8cm max height
%%\end{tikzpicture}
%
%\endofslide 
%%\endofsection
%\end{frame}
%%%%%%%%%%%%%%%


%%%%%%%%%%%%%%%
\begin{frame}
\frametitle{\texorpdfstring{\culine[\sectiontitleulinecolor]{\sectiontitle}}{\sectiontitle} }
\framesubtitle{Sähkövirrat pn-liitoksessa}
\appear{}
\seminormal%

\begin{itemize}[itemsep=1.5mm]
	\item Kohdistetaan liitokseen \CDAlert[fuchsia]{päästösuuntaus}\appear{, eli liitoksen yli muodostetaan positiivinen (sähkö)potentiaaliero}
	\appear{\xdef\loc{south east}%
\begin{tikzpicture}[remember picture,overlay] % anchor=south
    \node (A) [yshift=-0.25*\figYshift,xshift=1*\figXshift, anchor=\loc] at (current page.\loc) {\includegraphics[page=3,width=5cm]{Figures_booklet/Tikz_SOM_Bonding_figures_pn_junction}};%
\end{tikzpicture}%
}
	\item Päästösuuntaus heikentää tyhjennysalueen sähkökenttää\appear{, eli toisin sanoen p- ja n-puolien energiatasojen energiaeroa:}
\[
	\Delta E  \equal \appear{-e \cdot V}
\]
\appear{Tällöin elektronien on helpompi diffundoitua $p$-puolelle} \appear{ja aukkojen vastaavasti $n$-puolelle.} \appear{Tämä efekti \Alert[fuchsia]{kasvattaa molempia rekombinaatiovirtoja} Boltzmannin tekijällä:}
%$$
%e^{-\Delta E / k_{B} T}=e^{e \cdot V / k_{B} T}
%$$
%(We don't need Fermi-Dirac statistics, since most of the available states of diffusing electrons and holes are empty) 
\end{itemize}

% 6.5 cm = midway, 10 cm = 3/4 
\vspace{2.5mm}
\begin{minipage}{8.45cm}
	\begin{itemize}[itemsep=1.5mm]
\item[] $$
e^{-\Delta E / k_{B} T}  \equal \appear{e^{e \cdot V / k_{B} T}} \vspace{3mm}
$$
\appear{missä emme ole käyttäneet Fermin–Diracin statistiikkaa}\appear{, koska suurin osa tiloista joihin diffundoituvat tilat ovat siirtymässä ovat tyh\-jiä (ei tarvitse soveltaa kieltosääntöä}

\item Syntyy p-puolelta n-puolelle kulkeva virta 
\end{itemize}
\end{minipage}

\endofslide 
%\endofsection
\end{frame}
%%%%%%%%%%%%%%%

%%%%%%%%%%%%%%%%
%\begin{frame}
%\frametitle{\texorpdfstring{\culine[\sectiontitleulinecolor]{\sectiontitle}}{\sectiontitle} }
%
%\begin{itemize}
%	\item Now we apply forward bias, i.e. positive potential difference across the junction. A forward bias decreases the electric field in the junction region, as well as it decreases the differences between the energy levels on the $p$ - and $n$-sides by an amount of:
%$$
%\Delta E=-e \cdot V
%$$
%It becomes easier for electrons to diffuse to $p$-region and for the holes to $n$-region. This effect increases both recombination currents by the Boltzmann factor:
%$$
%e^{-\Delta E / k_{B} T}=e^{e \cdot V / k_{B} T}
%$$
%(We don't need Fermi-Dirac statistics, since most of the available states of diffusing electrons and holes are empty).
%\end{itemize}
%
%%\xdef\loc{south east}
%%\begin{tikzpicture}[remember picture,overlay] % anchor=south
%%    \node (A) [yshift=-0.25*\figYshift,xshift=\figXshift, anchor=\loc] at (current page.\loc) {\includegraphics[page=1,width=7cm]{Figures_booklet/SOM_10_Bonding_figures/SOM_Bonding_Figure_42.png}}; %scale=\figscale. % 8cm max height
%%\end{tikzpicture}
%
%\endofslide 
%%\endofsection
%\end{frame}
%%%%%%%%%%%%%%%%


%%%%%%%%%%%%%%%
\begin{frame}
\frametitle{\texorpdfstring{\culine[\sectiontitleulinecolor]{\sectiontitle}}{\sectiontitle} }
\framesubtitle{Sähkövirrat pn-liitoksessa}
\appear{}
\begin{itemize}
	\item Aukkojen kokonaisvirraksi voidaan nyt muodostaa:
\[
i_{p, t o t}   \equal \appear{i_{p, r}}  \minus \appear{| i_{p, g} | }  \equal  \appear{| i_{p, g} | }\appear{e^{e \cdot V / k_{B} T}}  \minus \appear{ | i_{p, g} |}  \equal  \appear{| i_{p, g} | }\appear{(e^{e \cdot V / k_{B} T}-1)}
\]
	\item Elektronien tapaus on vastaavanlainen\appear{, jolloin elektronien aiheuttamaksi kokonaisvirraksi saadaan}
\[
 I  \equal \appear{i_{p, t o t}}  \plus \appear{i_{n, t o t}} \equal \appear{I_{S}(T) ( e^{e \cdot V / k_{B} T}-1 )}
 \]
\appear{missä $I_{S}$ kutsutaan \Alert[red]{(estosuuntauksen)} \CDAlert[red]{saturaatiovirraksi}}\appear{, ja yhtälö tunnetaan nimellä \CDAlert[fuchsia]{Shockleyn diodiyhtälö} }

\item Tarkastelu voidaan toistaa estosuunnatulle liitokselle\appear{, jolloin diodiyhtälön huomataan pätevän molemmille tapauksille (positiivinen/negatiivinen ulkoinen kenttä)}

\item Saatu virta–jännite-ominaiskäyrä (I–V käyrä) esittää tyypillisen \CDAlert[fuchsia]{dioditasasuuntaajan} toiminnan\appear{, ja kuinka $I=I(V)$ käyttäytyy}%\appear{, \Alert[fuchsia]{dramaattisesti eroaa Ohmin laista}}
\end{itemize}

\endofslide 
%\endofsection
\end{frame}
%%%%%%%%%%%%%%%

%%%%%%%%%%%%%%%
\begin{frame}
\frametitle{\texorpdfstring{\culine[\sectiontitleulinecolor]{\sectiontitle}}{\sectiontitle} }
\framesubtitle{pn-liitoksen I–V käyrä}
\appear{}
\begin{itemize}
	\item Hahmottelemalla Shockleyn diodiyhtälön ennustaman virran $I(V)$
\[
I(V)  \equal  \appear{I_{S}(T)} \appear{( e^{e \cdot V / k_{B} T}-1 )}
\]
\appear{ulkoisen jännitteen $V$ funktiona}\appear{, huomaamme dramaattisen \CDAlert[fuchsia]{eron} \CDAlert[fuchsia]{Ohmin lakiin}}
\appear{\xdef\loc{south east}%
\begin{tikzpicture}[remember picture,overlay] % anchor=south
    \node (A) [yshift=-0.4*\figYshift,xshift=\figXshift, anchor=\loc] at (current page.\loc) {\includegraphics[page=1,width=7cm]{Figures_booklet/SOM_Tipler_Solid_state_physics_figures/SOM_Tipler_SSP_figure_46}};%scale=\figscale. % 8cm max height
\end{tikzpicture}\footnote{\TiplerCR}}
\end{itemize}

\vspace{2.5mm}
\begin{minipage}{6.7cm}
	\begin{itemize}%[itemsep=1.5mm]
\item Tämä on yksi diodin perusominaisuuksista\appear{, ja sitä on hankala saada aikaiseksi perinteisten komponenttien avulla} \appear{(vastukset}\appear{, kapasitaattorit} \appear{ja käämit)}

\item Todella suurilla estosuuntausjännitteillä \appear{(\CDAlert[blue]{läpilyöntijännite} = raja-arvo)}\appear{, pn-liitos alkaa päästää virtaa läpi myös toiseen suuntaan}
\end{itemize}
\end{minipage}

\endofslide 
%\endofsection
\end{frame}
%%%%%%%%%%%%%%%

%%%%%%%%%%%%%%%
\begin{frame}
\frametitle{\texorpdfstring{\culine[\sectiontitleulinecolor]{\sectiontitle}}{\sectiontitle} }
\framesubtitle{Valodiodi (LED)}
\appear{}
\seminormal
\begin{itemize}[itemsep=1.7mm]
	\item \CDAlert[red]{Valodiodi eli ledi (LED) on pn-liitos joka emittoi/säteilee valoa}
	
	\item Kun liitos on päästösuunnattu\appear{, monet aukot työntyvät niiden $\mathrm{p}$-alueelta liitoskohtaan (tyhjennysalueelle)}\appear{, samoin monet elektronit $n$-alueelta  }
	\implies{Nämä \Alert[red]{elektronit rekombinoituvat tyhjennysalueella aukkojen} kanssa}

\item Rekombinaation aikana\appear{ elektronit voivat emittoida fotonin}\appear{, jonka energia vastaa suurinpiirtein vyörakenteen energiarakoa} 

\end{itemize}

\vspace{1.7mm}
\begin{minipage}{8.7cm}
	\begin{itemize}[itemsep=1.7mm]
	\item Tätä energiaa \appear{(eli säteilevän valon väriä)} \appear{voidaan vaihdella muuttamalla pn-liitoksen materiaaleja, ja siten energiaraon suuruutta (ks.~edellinen kalvo) }

\item Ledejä käytetään nykyisin paljon valaistuksessa\appear{, digitaalisissa näytöissä,} \appear{jne.}

\item  $1$\;eV:n omaavan energiaraon (valo)diodi tarvitsee $1$\;V päästösuuntauksen syttyäkseen
\end{itemize}
\end{minipage}

\xdef\loc{south east}%
\begin{tikzpicture}[remember picture,overlay] % anchor=south
    \node (A) [yshift=-0.25*\figYshift,xshift=1*\figXshift, anchor=\loc] at (current page.\loc) {\includegraphics[page=5,width=5cm]{Figures_booklet/Tikz_SOM_Bonding_figures_pn_junction}};%
\end{tikzpicture}%

\endofslide 
%\endofsection
\end{frame}
%%%%%%%%%%%%%%%

%%%%%%%%%%%%%%%
\begin{frame}
\frametitle{\texorpdfstring{\culine[\sectiontitleulinecolor]{\sectiontitle}}{\sectiontitle} }
\framesubtitle{Valoilmaisin (valodiodi)}
\appear{}
% photodiode
\seminormal

%A Photodiode
\begin{itemize}[itemsep=2.2mm]
	\item Tarkastellaan suuntaamatonta \appear{(biasoimatonta)} \appear{diodia jonka energiarako on $1$\,eV}\appear{, ja jota valaistaan fotoneilla joiden energia on hieman suurempi kuin $1$\;eV} 
	\item \Alert[red]{Liitos absorboi fotoneita kaikilla alueilla nostaen elektroneita johtavuusvyölle} (ja aukkoja valenssivyölle) %A conduction electron in the p-side (where it is a minor carrier) is eager to slide down to the lower conduction band in the entype. Similarly holes created in either the depletion zone or $\mathrm{n}$-type valence zone would happily float up to the p-type valence. Thus, as long as the photons arrive from an external source, the component acts like a battery,
\end{itemize}

\vspace{2.1mm}
\begin{minipage}{7.8cm}
	\begin{itemize}[itemsep=2.2mm]
	%\item[] a hole in the valence band. 
	
	\item Johtavuus\CDAlert[red]{elektroni} p-puolella \appear{(missä se on vähemmistövarauksenkuljettaja)} \appear{\CDAlert[red]{laskeutuu} mielellään n-tyypin alemmalle johtavuusvyölle} 
	
	\item Vastaavasti \CDAlert[blue]{aukot} joita syntyy joko tyhjennysalueen tai n-puolen valenssivyöllä\appear{, mielellään \CDAlert[blue]{nousisivat} p-puolelle}
	
	\item Niin kauan kuin fotoneita saapuu liitokseen\appear{, komponentti toimii patterin tavoin } \appear{(\Alert[fuchsia]{eli synnyttää jännitteen}} \appear{\Alert[fuchsia]{mikä ylläpitää virtaa})}
	\end{itemize}
\end{minipage}

\xdef\loc{south east}
\begin{tikzpicture}[remember picture,overlay] % anchor=south
    \node (A) [yshift=-0.25*\figYshift,xshift=\figXshift, anchor=\loc] at (current page.\loc) {\includegraphics[page=1,width=5.7cm]{Figures_booklet/SOM_10_Bonding_figures/SOM_Bonding_Figure_43.png}}; %scale=\figscale. % 8cm max height
\end{tikzpicture}
\footnote{\HarrisCR[1-]}

\endofslide 
%\endofsection
\end{frame}
%%%%%%%%%%%%%%%

%%%%%%%%%%%%%%%
\begin{frame}
\frametitle{\texorpdfstring{\culine[\sectiontitleulinecolor]{\sectiontitle}}{\sectiontitle} }
\framesubtitle{Valosähkökenno}
\appear{}
\seminormal
\begin{itemize}[itemsep=1.8mm]
	\item Jos energiarako on $1$\;eV\appear{, niin fotonin energian täytyy olla ainakin yhtäsuuri kuin energiaraon:} \vspace{-2mm}
\[
\lambda \equal \frac{h c}{E_{rako}}  \equal \frac{1240\;eV \cdot nm}{1\; eV} \equal \appear{1240\;nm} \vspace{2mm}
\]
\item Jos tällainen laite kytketään virtapiiriin\appear{, se synnyttää sähkömotorisen voiman (emf) ja toimii virtalähteenä }
\implies{ Tällaista laitetta kutsutaan usein \CDAlert[fuchsia]{aurinkokennoksi tai valosähkökennoksi} } 

\item Laite ei tarvitse välttämättä auringonvaloa\appear{, koska riittää että fotonien energia ylittää energiaraon}

\item Nyt on tarkasteltu valosähkökennon toimintaperiaatetta\appear{, mikä muuntaa tulevan valon energiaksi}\appear{ mahdollistaen \CDAlert[fuchsia]{aurinkoenergian tuotannon}} %We may also add layers of intrinsic semiconductor between the n-type and $p$-type regions, widening the depletion zone where the electron-hole pairs are sorted quickly and contributing to improved efficiengy in high-frequency applications

\item Sama fysiikka pätee myös muihin valoherkkiin komponentteihin \appear{(eli \Alert[red]{valodiodeihin}):} \appear{varauskytkettyihin piireihin (CCD) perustuviin kamerakennoihin}\appear{, komplementaarisiin metallioksidipuolijohdedetektoreihin (\CDAlert[fuchsia]{CMOS})}\appear{, ja digitaalisiin kameroihin}
\end{itemize}

\endofslide 
%\endofsection
\end{frame}
%%%%%%%%%%%%%%%

%%%%%%%%%%%%%%%
\begin{frame}
\frametitle{\texorpdfstring{\culine[\sectiontitleulinecolor]{\sectiontitle}}{\sectiontitle} }
\framesubtitle{Diodilaseri}
%The Diode Laser
\appear{}
\seminormal
\begin{itemize}[itemsep=1.55mm]
	\item Kuten ledeissä\appear{, elektroni–aukkoparin rekombinaatiossa vapautuva energia} \appear{toimii \CDAlert[fuchsia]{diodilaserien} valonlähteenä}
	
	\item Suuret elektronien ja aukkojen konsentraatiot pn-liitoksessa \appear{mahdollistavat  \CDAlert[fuchsia]{populaatioinversion}}\appear{, mikä tarvitaan aineen laserointiin}

	\item Valon vahvistamiseksi tarvitaan myös optinen \CDAlert[fuchsia]{resonaattori} eli kaviteetti, mutta puolijohteiden suurten taitekertoimien takia  \appear{(esim. piillä $n\,{\sim}\,3,5$),}%due to the high refractive indices of many semiconductors (specifically for Si, $n \sim 3.7$), semiconductor–air interfaces often act as good enough cavity mirrors
\end{itemize}

%\vspace{1.7mm}
\begin{minipage}{8.2cm}
\begin{itemize}[itemsep=1.55mm]
\item[] puolijohde–ilma rajapinnat toimivat usein jo riittävän hyvinä resonaattorien peileinä

\item \CDAlert[red]{Epäsuoran energiaraon piitä} ei voi hyödyntää diodilaserin valmistukseen\appear{, koska se luovuttaa rekombinaatioenergiansa mieluummin muilla tavoilla kuin emittoimalla fotoneja}

\item \CDAlert[blue]{GaAs} on III–V puolijohde \appear{joka puolestaan on hyvä materiaali laserdiodin valmistukseen koska sillä on \CDAlert[blue]{suora energiarako} }

	\end{itemize}
\end{minipage}

\xdef\loc{south east}
\begin{tikzpicture}[remember picture,overlay] % anchor=south
    \node (A) [yshift=-0.4*\figYshift,xshift=\figXshift, anchor=\loc] at (current page.\loc) {\includegraphics[page=1,width=5.2cm,trim={0 2cm 0 0},clip]{Figures_booklet/SOM_Tipler_Solid_state_physics_figures/SOM_Tipler_SSP_figure_52.png}}; %scale=\figscale. % 8cm max height
\end{tikzpicture}
\footnote{\TiplerCR[1-]}

\endofslide 
%\endofsection
\end{frame}
%%%%%%%%%%%%%%%

%%%%%%%%%%%%%%%
\begin{frame}
\frametitle{\texorpdfstring{\culine[\sectiontitleulinecolor]{\sectiontitle}}{\sectiontitle} }
\framesubtitle{Transistori}
% The Transistor
\appear{}
\begin{itemize}[itemsep=1.6mm]
	\item Yksinkertaisimmillaan transistorit koostuvat kolmesta puolijohdealueesta\appear{, muodostaen joko \CDAlert[fuchsia]{npn-}} \appear{tai pnp-liitoksen (kuva)}
	\appear{\xdef\loc{south east}%
\begin{tikzpicture}[remember picture,overlay] % anchor=south
    \node (A) [yshift=-0.3*\figYshift,xshift=1*\figXshift, anchor=\loc] at (current page.\loc) {\includegraphics[page=1,width=13cm]{Figures_booklet/SOM_Tipler_Solid_state_physics_figures/SOM_Tipler_SSP_figure_53.png}};%scale=\figscale. % 8cm max height
\end{tikzpicture}\footnote{\TiplerCRshort}}
%\item Here npn is discussed

\item \Alert[red]{p-tyypin alue keskellä on  \CDAlert[red]{kanta}}\appear{, ja \Alert[blue]{n-tyypin alueet tunnetaan nimillä \CDAlert[blue]{emitteri}} \appear{ja \CDAlert[blue]{kollektori}}}

\item Emitteri and kollektori ovat erikokoisia ja niiden seostusasteet ovat erilaiset\appear{, koska niillä on erilaiset funktiot}
\end{itemize}

\endofslide 
%\endofsection
\end{frame}
%%%%%%%%%%%%%%%

%%%%%%%%%%%%%%%
\begin{frame}
\frametitle{\texorpdfstring{\culine[\sectiontitleulinecolor]{\sectiontitle}}{\sectiontitle} }
\framesubtitle{npn-liitoksinen transistori}
\appear{}
\begin{itemize}[itemsep=1.7mm]
	\item \CDAlert[blue]{Päästösuunnattu emitteri} diodi \appear{mahdollistaa johtavuuselektronien virtaamisen emitteristä} \appear{kantaan} 

\item Ohut ja  \CDAlert[red]{kevyesti seostettu kanta} on p-tyyppinen 
 \implies{e$^-$:t vähemmistövarauksenkuljettajina \appear{rekombinoituvat aukkojen kanssa harvoin}}
\end{itemize}
\vspace{3mm}
\begin{minipage}{5cm}
\begin{itemize}[itemsep=1.7mm]
\item Kollektori on n-tyypin puolijohde\appear{, ja \Alert[blue]{kollektori–kanta diodi on estosuunnattu} }

\item Pieni osa elektroneista rekombinoituu aukkojen kanssa kol\-lek\-to\-rissa\appear{, kun suurin osa laskeutuu kollektorin johta\-vuus\-vyölle}\appear{, synnyttäen \Alert[fuchsia]{virran jonka suuruutta voidaan} \CDAlert[fuchsia]{säätää kannalla}}
	\end{itemize}
\end{minipage}

\xdef\loc{south east}
\begin{tikzpicture}[remember picture,overlay] % anchor=south
    \node (A) [yshift=-0.25*\figYshift,xshift=\figXshift, anchor=\loc] at (current page.\loc) {\includegraphics[page=1,width=8.5cm]{Figures_booklet/SOM_10_Bonding_figures/SOM_Bonding_Figure_44.png}}; %scale=\figscale. % 8cm max height
\end{tikzpicture}
\footnote{\HarrisCR[1-]}
\endofslide 
%\endofsection
\end{frame}
%%%%%%%%%%%%%%%

%%%%%%%%%%%%%%%
\begin{frame}
\frametitle{\texorpdfstring{\culine[\sectiontitleulinecolor]{\sectiontitle}}{\sectiontitle} }
\framesubtitle{Transistorivahvistin}
\appear{}
\begin{itemize}
	\item Kollektorista ulos virtaavien elektronien osuus \appear{(palaavat toki suljetussa piirissä lopulta takaisin emitteriin)} \appear{on \CDAlert[red]{verrannollinen}} %and typically 100 times greater than that which flows out the base (returning through the input circuit).
\end{itemize}

\begin{minipage}{6.75cm}
\begin{itemize}%[itemsep=1.7mm]
\item[]  \Alert[red]{signaaliin} \appear{\Alert[red]{ja tyypillisesti 100-kertainen kantaan tulevaan signaaliin nähden}}% 
\appear{\xdef\loc{south east}%
\begin{tikzpicture}[remember picture,overlay] % anchor=south
    \node (A) [yshift=-0.35*\figYshift,xshift=\figXshift, anchor=\loc] at (current page.\loc) {\includegraphics[page=1,width=6.8cm]{Figures_booklet/SOM_10_Bonding_figures/SOM_Bonding_Figure_45.png}}; %scale=\figscale. % 8cm max height
\end{tikzpicture}\footnote{\HarrisCR}%
}%
\item Kuvassa~45 \appear{sisääntulo(signaali) aiheuttaa pieniä jännitevaihteluita} \appear{emitteri–kannan potentiaali-eroon}\appear{, mikä johtaa vastaavaan vaihteluun} \appear{kannan ja emitterin läpimenovirrassa} \appear{(virtaavien elektronien nopeus vaihtelee).} \appear{Pieni osa tästä virrasta poistuu kannasta} \appear{ja kiertää moduloituna signaalina (emitteri–kanta) virtapiirissä}
	\end{itemize}
\end{minipage}

\endofslide 
%\endofsection
\end{frame}
%%%%%%%%%%%%%%%

%%%%%%%%%%%%%%%
\begin{frame}
\frametitle{\texorpdfstring{\culine[\sectiontitleulinecolor]{\sectiontitle}}{\sectiontitle} }
\framesubtitle{Kanavatransistori (FET)}
\appear{}
\semismall
\begin{itemize}[itemsep=1.4mm]
	\item \Alert[red]{Kanavatransistori (FET)} on yleinen kytkintyyppi \CDAlert[fuchsia]{integroiduissa piireissä}
	
	\item Kanavatransistori voidaan muodostaa p-tyypin piilevylle \appear{seostamalla levyn pinnalle kaksi n-tyypin aluetta}\appear{, joita kutsutaan '\CDAlert[blue]{lähteeksi}' ja '\CDAlert[tuniblue]{nieluksi}'}
	
	\item  Keskelle jääneen p-tyypin päälle laitetaan elektrodi \appear{jota kutsutaan '\CDAlert[red]{portiksi}'} \appear{ja joka on erotettu levystä}\appear{ eristeenä toimivalla piidioksidikerroksella ($\mathrm{SiO}_{2}$)} 	
		
\end{itemize}

\vspace{1.4mm}
\begin{minipage}{7.3cm}
\begin{itemize}[itemsep=1.4mm]
	\item Kun portti on sähköisesti neutraali\appear{ vain pieni virta menee lähteen ja nielun läpi}\appear{, koska toinen näistä pn-liitoksista on estosuunnattu } 
	
	\item Positiivisesti varautuneessa \Alert[red]{pienessä} portissa\appear{ jo pieni varaus }\appear{synnyttää suu\-ren kentän p-alueelle}\appear{, mahdollistaen  \CDAlert[red]{virran liikkumisen koko npn-liitoksen läpi}  }
	
	\item Kanavatransistorit ovat pieniä \appear{(${\sim}\,1\;\mu m$)} \appear{integ\-roituja transistoreita}
	
	\item Metallioksidipuolijohdetransistori \\ 
	         (\Alert[red]{MOS-transistori}) %\Alert[red]{MOSFET}{https://en.wikipedia.org/wiki/MOSFET\#Importance} = Metal–oxide–semiconductor FET
	\end{itemize}
\end{minipage}

\xdef\loc{south east}
\begin{tikzpicture}[remember picture,overlay] % anchor=south
    \node (A) [yshift=-0.25*\figYshift,xshift=\figXshift, anchor=\loc] at (current page.\loc) {\includegraphics[page=1,width=6.5cm]{Figures_booklet/SOM_10_Bonding_figures/SOM_Bonding_Figure_46.png}}; %scale=\figscale. % 8cm max height
\end{tikzpicture}
\footnote{\HarrisCR}

%\endofslide 
\endofsection
\end{frame}
%%%%%%%%%%%%%%%

%%%%%%%%%%%%%%%%
%\begin{frame}
%\frametitle{\texorpdfstring{\culine[\sectiontitleulinecolor]{\sectiontitle}}{\sectiontitle} }
%
%\begin{itemize}
%	\item A single-electron transistor (SET) is a sensitive electronic device based on the Coulomb blockade effect
%
%\item In this device the electrons flow through a tunnel junction between source/drain to a quantum dot (conductive island)
%\end{itemize}
%
%%\xdef\loc{south east}
%%\begin{tikzpicture}[remember picture,overlay] % anchor=south
%%    \node (A) [yshift=-0.25*\figYshift,xshift=\figXshift, anchor=\loc] at (current page.\loc) {\includegraphics[page=1,width=7cm]{Figures_booklet/SOM_10_Bonding_figures/SOM_Bonding_Figure_42.png}}; %scale=\figscale. % 8cm max height
%%\end{tikzpicture}
%
%\endofslide 
%%\endofsection
%\end{frame}
%%%%%%%%%%%%%%%%
